\textbf{Open-set Action Recognition.}
In open-set action recognition (OSAR), the model is required to recognize the known action samples from training classes and reject the unknown samples that are out of training classes. 
Compared with images, video actions are more challenging to be recognized in an OSR setting than images due to the uncertain temporal dynamics, and static bias of human actions \cite{bao2021evidential}. Our \model~generalizes well to unknown classes that are out of training classes and outperforms the existing method \cite{bao2021evidential} without any model calibration.


We use the ViT-H/16 model of \model~as a backbone, and finetune it with a simple linear classification head on UCF-101 \cite{soomro2012dataset} training set. To enable \model~to ``know unknown", we follow the method DEAR proposed in \cite{bao2021evidential} and formulate it as an uncertainty estimation problem by leveraging evidential deep learning (EDL), which provides a way to jointly formulate the multiclass classification and uncertainty modeling. Specifically, given a video as input, the Evidential Neural Network (ENN) head on top of a \model~backbone predicts the class-wise evidence, which formulates a Dirichlet distribution so that the multi-class probabilities and predictive uncertainty of the input can be determined. During the open-set inference, high-uncertainty videos can be regarded as unknown actions, while low-uncertainty videos are classified by the learned categorical probabilities. 

\model~can not only recognize known action classes accurately but also identify the unknown. Table \ref{osar} reports the results of both closed-set (Closed Set Accuracy) and open-set (Open Set AUC) performance of \model~and other baselines. It shows that our \model~consistently and significantly outperforms other baselines on both two open-set datasets, where unknown samples are from HMDB-51~\cite{hmdb51} and MiT-v2 \cite{monfort2021multi}, respectively.
